\documentclass[aspectratio=169,11pt]{beamer}
% Shared Beamer style for practice contest solution walkthroughs.
\usepackage[utf8]{inputenc}
\usepackage[T1]{fontenc}
\usepackage{lmodern}
\usepackage{amsmath,amssymb}
\usepackage{xcolor}
\usepackage{hyperref}
\usepackage{listings}

\definecolor{CPNavy}{HTML}{172033}
\definecolor{CPBlue}{HTML}{2563EB}
\definecolor{CPGreen}{HTML}{16A34A}
\definecolor{CPOrange}{HTML}{F97316}
\definecolor{CPGray}{HTML}{64748B}
\definecolor{CPLight}{HTML}{F8FAFC}
\definecolor{CPCodeBg}{HTML}{F1F5F9}
\definecolor{CPCodeRule}{HTML}{CBD5E1}

\usetheme{default}
\usefonttheme{professionalfonts}
\setbeamertemplate{navigation symbols}{}
\setbeamersize{text margin left=0.65cm,text margin right=0.65cm}
\setbeamertemplate{itemize item}{\color{CPBlue}$\blacktriangleright$}
\setbeamertemplate{itemize subitem}{\color{CPGray}$\triangleright$}

\setbeamercolor{normal text}{fg=CPNavy,bg=white}
\setbeamercolor{frametitle}{fg=white,bg=CPNavy}
\setbeamercolor{title}{fg=white,bg=CPNavy}
\setbeamercolor{subtitle}{fg=CPLight,bg=CPNavy}
\hypersetup{colorlinks=true,urlcolor=CPBlue}

\setbeamertemplate{frametitle}{%
  \nointerlineskip%
  \begin{beamercolorbox}[wd=\paperwidth,ht=0.88cm,dp=0.22cm,leftskip=0.55cm,rightskip=0.35cm]{frametitle}
    \usebeamerfont{frametitle}\insertframetitle\hfill\small\insertframenumber/\inserttotalframenumber
  \end{beamercolorbox}%
}

\setbeamertemplate{footline}{%
  \leavevmode%
  \hbox{%
    \begin{beamercolorbox}[wd=.58\paperwidth,ht=2.4ex,dp=1ex,leftskip=.5cm]{author in head/foot}%
      \color{CPGray}\insertshorttitle
    \end{beamercolorbox}%
    \begin{beamercolorbox}[wd=.42\paperwidth,ht=2.4ex,dp=1ex,rightskip=.5cm plus1fil]{date in head/foot}%
      \hfill\color{CPGray}\insertshortauthor
    \end{beamercolorbox}%
  }%
}

\setbeamertemplate{title page}{%
  \vspace*{1.25cm}
  \begin{beamercolorbox}[wd=\paperwidth,sep=0.65cm,leftskip=0.15cm,rightskip=0.15cm]{title}
    {\color{white}\Huge\bfseries\inserttitle\par}
    \vspace{0.35cm}
    {\color{CPLight}\Large\insertsubtitle\par}
    \vspace{0.6cm}
    {\color{CPLight}\large\insertauthor\par}
  \end{beamercolorbox}
  \vfill
}

\newcommand{\sectiontag}[1]{\textbf{\color{CPBlue}#1}}
\newenvironment{tightitemize}{\begin{itemize}\small\setlength{\itemsep}{0.18cm}}{\end{itemize}}
\lstdefinestyle{cpcode}{
  language=Python,
  basicstyle=\ttfamily\tiny,
  keywordstyle=\color{CPBlue}\bfseries,
  commentstyle=\color{CPGray}\itshape,
  stringstyle=\color{CPGreen},
  backgroundcolor=\color{CPCodeBg},
  frame=single,
  framerule=0.25pt,
  rulecolor=\color{CPCodeRule},
  showstringspaces=false,
  columns=fullflexible,
  keepspaces=true,
  breaklines=true,
  breakatwhitespace=false,
  tabsize=4,
  xleftmargin=0.08cm,
  xrightmargin=0.08cm,
  aboveskip=0.08cm,
  belowskip=0.08cm
}


\title[imperfectgps]{Imperfect GPS}
\subtitle{Day 5 solution walkthrough}
\author[Solution Deck]{Competitive Programming Bootcamp}
\date{}

\begin{document}

\begin{frame}[plain]
  \titlepage
\end{frame}

% BEGIN GENERATED STATEMENT INTERPRETATION
\begin{frame}{Interpret the Word Problem}
\small
\sectiontag{Statement excerpts:}
\begin{tightitemize}
  \item \textit{``records its position periodically''}
  \item \textit{``GPS can underestimate the total distance''}
\end{tightitemize}

\vspace{0.18cm}
\sectiontag{Programming interpretation:}
\begin{tightitemize}
  \item Compute the true path length from every original segment.
  \item Compute the GPS path from sampled positions at fixed time intervals plus the final point.
  \item The output is the percentage of true distance lost by sampling.
\end{tightitemize}
\end{frame}
% END GENERATED STATEMENT INTERPRETATION







\begin{frame}{Problem Model}
\small
\sectiontag{Textbook connection:} CP4 7.2.1 Points

\vspace{0.25cm}
\sectiontag{Core technique:} Geometry simulation with interpolation

\vspace{0.25cm}
\begin{tightitemize}
  \item The actual path is the polyline through all recorded runner positions.
  \item The GPS samples the position every fixed interval and at the final time.
  \item The answer is the percentage distance lost by GPS sampling.
\end{tightitemize}
\end{frame}

\begin{frame}{How to Recognize the Approach}
\begin{tightitemize}
  \item Actual distance is a sum of Euclidean segment lengths.
  \item GPS samples may fall between known points.
  \item Constant speed on each segment allows linear interpolation.
\end{tightitemize}
\end{frame}

\begin{frame}{Algorithm}
\begin{tightitemize}
  \item Compute actual\_distance across consecutive input points.
  \item Start previous GPS point at the initial position.
  \item For each GPS time before the final time, find its containing segment.
  \item Interpolate x and y using the time ratio.
  \item Add GPS segment distances, then add the final position and compute percentage lost.
\end{tightitemize}
\end{frame}

\begin{frame}{Walkthrough of the Key State}
\begin{tightitemize}
  \item ratio = (gps\_time - time1) / (time2 - time1).
  \item current = point1 + ratio * (point2 - point1).
  \item math.hypot computes stable Euclidean distances.
\end{tightitemize}
\end{frame}

\begin{frame}{Why the Algorithm Is Correct}
\begin{tightitemize}
  \item Actual distance sums exactly the straight segments described by the input path.
  \item Interpolation gives the runner's exact position at every GPS sample time under constant segment speed.
  \item The GPS distance is the polyline through those sampled positions, matching the device model.
\end{tightitemize}
\end{frame}

\begin{frame}{Complexity and Implementation Notes}
\small
\sectiontag{Complexity:} O(N + T/interval) time and O(N) memory.

\vspace{0.3cm}
\begin{tightitemize}
  \item The final GPS point is always included even when the final time is not a multiple of the interval.
  \item Use floating-point output; Kattis accepts tolerance.
\end{tightitemize}
\end{frame}



% BEGIN GENERATED CODE WALKTHROUGH
\begin{frame}[fragile]{Code: Read Timed Points}
\scriptsize
\sectiontag{Source lines:} \texttt{imperfectgps.py:46--64}

\vspace{0.08cm}
\begin{columns}[T,totalwidth=\textwidth]
\begin{column}{0.58\textwidth}
\begin{lstlisting}[style=cpcode]
import sys
from math import hypot

def distance(x1, y1, x2, y2):
    return hypot(x2 - x1, y2 - y1)

def main():
    input = sys.stdin.buffer.readline

    n, interval = map(int, input().split())

    points = []

    for _ in range(n):
        x, y, time = map(int, input().split())
        points.append((x, y, time))
\end{lstlisting}
\end{column}
\begin{column}{0.38\textwidth}
\sectiontag{What this chunk does}
\begin{tightitemize}
  \item hypot computes Euclidean segment lengths accurately.
  \item Each input point stores x, y, and timestamp.
  \item The interval controls when GPS readings are simulated.
\end{tightitemize}
\end{column}
\end{columns}
\end{frame}

\begin{frame}[fragile]{Code: Actual Distance}
\scriptsize
\sectiontag{Source lines:} \texttt{imperfectgps.py:66--76}

\vspace{0.08cm}
\begin{columns}[T,totalwidth=\textwidth]
\begin{column}{0.58\textwidth}
\begin{lstlisting}[style=cpcode]
actual_distance = 0.0

for i in range(1, n):
    x1, y1, _ = points[i - 1]
    x2, y2, _ = points[i]

    actual_distance += distance(x1, y1, x2, y2)
\end{lstlisting}
\end{column}
\begin{column}{0.38\textwidth}
\sectiontag{What this chunk does}
\begin{tightitemize}
  \item The real route follows every consecutive input point.
  \item Each segment length is added to actual\_distance.
  \item This is the denominator for the final percentage loss.
\end{tightitemize}
\end{column}
\end{columns}
\end{frame}

\begin{frame}[fragile]{Code: Initialize GPS Sampling}
\scriptsize
\sectiontag{Source lines:} \texttt{imperfectgps.py:78--95}

\vspace{0.08cm}
\begin{columns}[T,totalwidth=\textwidth]
\begin{column}{0.58\textwidth}
\begin{lstlisting}[style=cpcode]
previous_x = points[0][0]
previous_y = points[0][1]

gps_distance = 0.0

gps_time = interval

segment = 0

final_time = points[-1][2]
\end{lstlisting}
\end{column}
\begin{column}{0.38\textwidth}
\sectiontag{What this chunk does}
\begin{tightitemize}
  \item The GPS always records the starting position.
  \item previous\_x and previous\_y track the last recorded GPS point.
  \item segment identifies which true path segment contains the next GPS timestamp.
\end{tightitemize}
\end{column}
\end{columns}
\end{frame}

\begin{frame}[fragile]{Code: Interpolate Regular Samples}
\scriptsize
\sectiontag{Source lines:} \texttt{imperfectgps.py:97--135}

\vspace{0.08cm}
\begin{columns}[T,totalwidth=\textwidth]
\begin{column}{0.58\textwidth}
\begin{lstlisting}[style=cpcode]
while gps_time < final_time:

    while points[segment + 1][2] < gps_time:
        segment += 1

    x1, y1, time1 = points[segment]
    x2, y2, time2 = points[segment + 1]

    ratio = (gps_time - time1) / (time2 - time1)

    current_x = x1 + ratio * (x2 - x1)
    current_y = y1 + ratio * (y2 - y1)

    gps_distance += distance(previous_x, previous_y, current_x, current_y)

    previous_x = current_x
    previous_y = current_y

    gps_time += interval
\end{lstlisting}
\end{column}
\begin{column}{0.38\textwidth}
\sectiontag{What this chunk does}
\begin{tightitemize}
  \item The loop processes each scheduled GPS time before the finish.
  \item The segment pointer advances until the timestamp lies in the correct true segment.
  \item Linear interpolation finds the runner's position at that time before adding GPS distance.
\end{tightitemize}
\end{column}
\end{columns}
\end{frame}

\begin{frame}[fragile]{Code: Finalize Percentage Loss}
\scriptsize
\sectiontag{Source lines:} \texttt{imperfectgps.py:137--159}

\vspace{0.08cm}
\begin{columns}[T,totalwidth=\textwidth]
\begin{column}{0.58\textwidth}
\begin{lstlisting}[style=cpcode]
    final_x = points[-1][0]
    final_y = points[-1][1]

    gps_distance += distance(previous_x, previous_y, final_x, final_y)

    lost_distance = actual_distance - gps_distance

    percentage_lost = (lost_distance / actual_distance) * 100

    print(percentage_lost)

main()
\end{lstlisting}
\end{column}
\begin{column}{0.38\textwidth}
\sectiontag{What this chunk does}
\begin{tightitemize}
  \item The final position is always recorded, even off the regular interval.
  \item lost\_distance compares actual path length with sampled GPS path length.
  \item The percentage is printed as a floating-point value.
\end{tightitemize}
\end{column}
\end{columns}
\end{frame}
% END GENERATED CODE WALKTHROUGH

\begin{frame}{Source}
\small
\sectiontag{Solution file:} \texttt{practice\_contests/day\_5/imperfectgps.py}

\vspace{0.3cm}
\sectiontag{Problem:} \url{https://open.kattis.com/problems/imperfectgps}

\vspace{0.45cm}
Use this deck with the source code open beside it. The slides explain the reasoning; the code shows the exact input handling and output formatting.
\end{frame}

\end{document}
